% Chapter 10 -- Claim: The Yield Curve Predicts Future Interest Rates
% Requires: amsmath, amssymb, booktabs, array, tabularx, graphicx, tcolorbox (breakable, skins), hyperref or url.
% Defines the callout box only if the preamble has not.
\providecolor{fenavy}{HTML}{1B2A6B}
\providecolor{feshade}{HTML}{F3F5FA}
\makeatletter
\@ifundefined{callout}{%
  \newtcolorbox{callout}{breakable, enhanced,
    colback=white, colframe=gray!70, boxrule=0.5pt, arc=1mm,
    borderline west={2.5pt}{0pt}{gray!70},
    left=8pt, right=8pt, top=4pt, bottom=4pt}%
}{}
\makeatother

\chapter{Claim: The Yield Curve Predicts Future Interest Rates}
\label{ch:10}

When the yield curve\index{yield curve} slopes steeply upward, commentators read it as a forecast: markets expect interest rates to rise. When it inverts, they read it as a forecast of cuts and recession. Central banks extract ``market-implied'' policy paths from forward rates\index{forward rate}, and traders quote them as the market's view.

That reading rests on the oldest theory of the term structure, the \textbf{expectations hypothesis}\index{expectations hypothesis|textbf}: long rates are averages of expected future short rates, so forward rates are unbiased forecasts of where short rates are going. If it holds, a long bond and a sequence of short bonds earn the same expected return\index{return!expected}, and there is nothing to gain from choosing maturities.

The hypothesis is also one of the most thoroughly tested claims in finance. This chapter brings the forecasting tools of Part III to the bond market with that test. The answer matters for anyone who reads the yield curve, manages a bond portfolio or prices an interest-rate derivative.

The chapter is built around three papers on the yield curve. Fama and Bliss (1987) and Campbell and Shiller (1991) tested whether long rates are good forecasts of future short rates, as the expectations hypothesis says. Cochrane and Piazzesi (2005) showed that a single combination of forward rates predicts bond excess returns. The lab in Section~\ref{sec:10-lab} replicates these regressions on the Treasury curve since 1961 in order to check whether the curve forecasts future rates, or instead forecasts the extra return that long bonds pay.

By the end of the chapter you will be able to:

\begin{itemize}
  \item compute zero-coupon\index{yield curve!zero-coupon} yields, forward rates, holding-period returns\index{return!holding-period} and excess returns\index{return!excess} from bond prices;
  \item use duration\index{duration} to estimate a bond's price sensitivity to its yield;
  \item state the expectations hypothesis and define the term premium\index{term premium};
  \item explain the Fama--Bliss\index{Fama--Bliss regression}, Campbell--Shiller\index{Campbell--Shiller decomposition} and Cochrane--Piazzesi\index{Cochrane--Piazzesi factor} tests and why they ask the same question;
  \item summarize the evidence, including the small-sample and out-of-sample\index{in-sample versus out-of-sample} critiques;
  \item read a market-implied rate path as a forecast plus a premium, and judge a bond strategy built on carry and roll-down\index{roll-down}.
\end{itemize}

\section{Foundations: yields, forward rates and the term premium}
\label{sec:10-foundations}

\subsection{Zero-coupon yields}

Bond yields, the yield curve and the term spread\index{term spread} appear throughout finance. This chapter needs a precise version of them built from zero-coupon bonds. A zero-coupon bond pays \$1 at maturity $n$ and nothing before. Its price $P(n)$ defines its yield $y(n)$. With continuous compounding\index{compounding}, $P(n) = e^{-n\,y(n)}$, so the log price is
\begin{equation}
  p(n) = -n\,y(n)
  \label{eq:10-logprice}
\end{equation}
Coupon bonds are portfolios of zeros, so the \textbf{zero-coupon yield curve}\index{yield curve!zero-coupon|textbf} (yield against maturity) summarizes all Treasury prices. The Federal Reserve publishes a fitted curve back to 1961 (G\"urkaynak, Sack and Wright 2007).

\subsection{Forward rates}

Buying an $n$-year zero and selling an $(n-1)$-year zero locks in today a rate of return for the year from $n-1$ to $n$. That rate is the \textbf{forward rate}\index{forward rate|textbf} (Figure~\ref{fig:10-forward}):
\begin{equation}
  f(n) = p(n-1) - p(n) = n\,y(n) - (n-1)\,y(n-1)
  \label{eq:10-forward}
\end{equation}
A rising yield curve means forward rates above current yields.

\begin{figure}[htbp]
\centering
\scalebox{0.8}{%
\begin{tikzpicture}[x=3.2cm, y=1cm]
  % Bars: 1-year zero (0-1), 2-year zero (0-2), forward (1-2)
  \fill[gray!45] (0,0.6) rectangle (1,1.15);
  \node at (0.5,0.875) {1-year zero: 4\%};
  \fill[gray!22] (0,1.55) rectangle (2,2.1);
  \node at (1,1.825) {2-year zero: 5\% a year};
  \fill[fenavy] (1,2.5) rectangle (2,3.05);
  \node[white] at (1.5,2.775) {forward: 6\%};

  \node[anchor=north west, align=left] at (2.12,3.2)
    {$f(2) = 2 \times 5\% - 1 \times 4\% = 6\%$:\\the rate locked in today\\for the second year};

  % Time axis
  \draw[-stealth] (0,0) -- (2.35,0) node[anchor=west] {years};
  \foreach \wfT in {0,1,2} {
    \draw (\wfT,0.08) -- (\wfT,-0.08);
    \node[anchor=north] at (\wfT,-0.1) {\wfT};
  }
\end{tikzpicture}%
}
\caption{A forward rate from the worked example. A two-year yield of 5\% is equivalent to earning the one-year yield of 4\% in the first year and a rate locked in today, the forward rate of 6\%, in the second.}
\label{fig:10-forward}
\end{figure}

\subsection{Holding-period returns}

Buying an $n$-year bond and selling it a year later, when it is an $(n-1)$-year bond, earns the log return\index{return!log}
\begin{equation}
  r_{t+1}(n) = p_{t+1}(n-1) - p_t(n)
  \label{eq:10-hpr}
\end{equation}
As for stocks (Section~\ref{sec:01-rf}), the excess return $rx_{t+1}(n)$ subtracts the return on the safe alternative, here the one-year yield $y_t(1)$, which is the return on a one-year bond held to maturity.

\subsection{Duration}

The percentage price change of a bond for a small change in its yield is approximately $-D \times \Delta y$, where $D$ is the \textbf{duration}\index{duration|textbf}. For a zero, $D$ equals its maturity. Long bonds are more sensitive to rates, which is why their returns are more volatile.

\subsection{The expectations hypothesis (EH)}

\textbf{The expectations hypothesis (EH)}\index{expectations hypothesis|textbf} states that the forward rate equals the expected future one-year rate, up to a constant:
\begin{equation}
  f_t(n) = E_t\!\left[y_{t+n-1}(1)\right] + \text{constant}
  \label{eq:10-eh}
\end{equation}
Equivalently, expected excess returns on bonds of all maturities are constant over time. The \textbf{term premium}\index{term premium|textbf} is the gap between the forward rate and the expected future short rate, or between a long yield and the average expected short rate over its life. The EH says the term premium does not vary.

\textbf{Worked example.} The one-year zero yields 4\% and the two-year zero 5\% (continuously compounded).

\begin{itemize}
  \item Log prices: $p(1) = -0.04$, $p(2) = -0.10$. The one-year forward rate one year ahead is $f(2) = -0.04 + 0.10 = 6\%$.
  \item Under the EH, the market expects the one-year rate to be 6\% next year. If it is, the two-year bond, sold next year as a one-year bond, returns $-0.06 + 0.10 = 4\%$, the same as the one-year bond: excess return zero.
  \item If instead the one-year rate stays at 4\%, the two-year bond returns $-0.04 + 0.10 = 6\%$, an excess return of 2\%.
  \item A steep curve therefore forecasts either rising short rates (the EH) or high excess returns on long bonds (a term premium), or some mix (Figure~\ref{fig:10-readings}). The regressions in this chapter measure the mix.
\end{itemize}

\begin{figure}[htbp]
\centering
\scalebox{0.8}{%
\begin{tikzpicture}[
  trbox/.style={draw=fenavy, thick, rounded corners=2pt, align=center,
                minimum width=4.2cm, minimum height=1.15cm, inner sep=4pt},
  trarr/.style={-stealth, thick},
]
  \node[trbox, fill=fenavy!12] (top) at (0,0)
    {Steep yield curve:\\forward rates above short rates};
  \node[trbox] (eh) at (-3.4,-2.4)
    {short rates will rise\\(expectations hypothesis)};
  \node[trbox] (tp) at (3.4,-2.4)
    {long bonds will earn\\more (term premium)};
  \draw[trarr] (top.south) ++(-0.8,0) -- (eh.north);
  \draw[trarr] (top.south) ++(0.8,0) -- (tp.north);
  \node at (0,-3.7) {the regressions in this chapter measure the mix};
\end{tikzpicture}%
}
\caption{Two readings of a steep yield curve. Forward rates above short rates can mean that short rates will rise, as the expectations hypothesis says, or that long bonds will earn more than short ones, a term premium. The tests in this chapter measure how much of each.}
\label{fig:10-readings}
\end{figure}

\section{How it works}
\label{sec:10-how}

The expectations hypothesis can be tested in three ways that look different but ask the same question.

\begin{center}
\small
\begin{tabularx}{0.9\linewidth}{@{}>{\raggedright\arraybackslash}p{0.17\linewidth}>{\raggedright\arraybackslash}X>{\raggedright\arraybackslash}p{0.15\linewidth}>{\raggedright\arraybackslash}X@{}}
\toprule
Test & Regression & Expectations hypothesis & Alternative: time-varying term premium \\
\midrule
Fama--Bliss, excess returns & $rx(n)$ on $f(n) - y(1)$ & Slope 0 & Slope positive (1 if forward spreads are all premium) \\
Fama--Bliss, short rates & $y(1)$ in $n-1$ years minus $y(1)$ today, on $f(n) - y(1)$ & Slope 1 & Slope below 1 \\
Campbell--Shiller & Change in the long yield over a year, on the scaled yield spread & Slope 1 & Slope below 1, even negative \\
\bottomrule
\end{tabularx}
\end{center}

The first two are mirror images. The forward spread equals the expected change in short rates plus the expected excess return, so a spread that does not forecast rising rates must forecast excess returns. The Campbell--Shiller test asks the question through long yields: if the curve is steep because short rates are expected to rise, long yields should rise too, as the bond approaches the higher short rates. If they fall instead, the steepness was a premium.

\textbf{Why the premium might vary.} Bonds lose value when inflation\index{inflation} or rates rise unexpectedly. Investors demand compensation for that risk, and the compensation can change: higher in recessions, when investors are more risk-averse\index{risk aversion}, and when inflation is uncertain. Supply and demand matter too. When pension funds and central banks buy long bonds, they push the term premium down. Any of these makes the EH fail.

\textbf{Why the tests are hard.} Interest rates are highly persistent, and annual returns sampled monthly overlap by eleven months. Both inflate $t$-statistics\index{t-statistic} if handled carelessly (Chapter~\ref{ch:03}). The lab uses Newey--West\index{standard error!Newey--West (HAC)} errors with 18 lags and adds out-of-sample tests (Chapter~\ref{ch:08}).

\section{The claim}
\label{sec:10-claim}

\begin{callout}
\textbf{Claim (the expectations hypothesis).} Forward rates are unbiased forecasts of future short-term interest rates, up to a constant. A steep yield curve predicts rising short rates, not higher returns on long bonds.
\end{callout}

To test it we need to be precise on three points.

\begin{itemize}
  \item \textbf{Up to a constant.} The pure EH says forward rates equal expected short rates. Almost no one believes that, because long bonds have earned more than short ones on average. The testable version allows a constant term premium and denies only that it \emph{varies}.
  \item \textbf{Which horizon.} The tests use one-year holding periods, the standard in the literature. Fama and Bliss also examine longer horizons.
  \item \textbf{What would count against it.} The EH fails if forward spreads forecast excess bond returns (Fama--Bliss slope significantly above 0), if they forecast changes in short rates less than one for one, or if long yields fail to rise when the curve is steep (Campbell--Shiller slope significantly below 1). It fails convincingly if these patterns hold out of sample.
\end{itemize}

\section{The original evidence}
\label{sec:10-evidence}

\textbf{Fama and Bliss (1987).} Using discount bond prices constructed from CRSP\index{CRSP} Treasury data, Fama and Bliss found that the spread between the $n$-year forward rate and the one-year rate forecast one-year excess returns on $n$-year bonds, contradicting the EH. The same spread had little power to forecast changes in the one-year rate over the next year, but its power rose with the horizon: forward rates did contain information about short rates several years ahead, as interest rates mean-revert.

\textbf{Campbell and Shiller (1991).} Across maturities from a few months to ten years, the yield spread forecast long yields to move in the \emph{wrong} direction. When the curve was steep, long yields tended to fall over the following period, the opposite of what the EH requires, and the slope coefficients became more negative as maturity increased. A steep curve forecast high returns on long bonds, through rising prices.

\textbf{Cochrane and Piazzesi (2005).} Instead of one forward spread at a time, they regressed excess returns on all forward rates at once. A single factor, a tent-shaped combination of forward rates, forecast one-year excess returns on bonds of all maturities, with $R^2$\index{regression!R-squared@$R^2$} values up to 43\% in their sample. The factor was largely unrelated to the level, slope and curvature of the yield curve and rose in bad times, consistent with a risk premium\index{risk premium} that is high in recessions.

\textbf{The skeptics.} The evidence for the EH's failure is now widely accepted. The strength of the more ambitious claims is not. Bauer and Hamilton (2018) showed that standard tests suffer serious small-sample distortions with persistent predictors and overlapping returns\index{overlapping observations}. With more robust tests, the evidence that variables beyond level, slope and curvature forecast bond returns, including the Cochrane--Piazzesi factor, became much weaker. Thornton and Valente (2012) found that forecasts from forward rates gave little economic value out of sample to an investor choosing between bonds.

\section{Lab~\thechapter: The Treasury curve since 1961}
\label{sec:10-lab}

\begin{callout}
\textbf{Lab~\ref{ch:10} at a glance.} \emph{Claim on trial\index{claims-on-trial method}:} forward rates are unbiased forecasts of future short rates. \emph{Data:} the Federal Reserve's zero-coupon Treasury yield curve (G\"urkaynak, Sack and Wright 2007), month-end yields for maturities of one to ten years; free and updated regularly. \emph{Code:} runs as a script (\texttt{python lab16.py}) or as a notebook with a prediction prompt before each step.
\end{callout}

\textbf{Part A, the classic tests.} Fama--Bliss regressions of one-year excess returns on forward spreads for maturities of 2 to 10 years, the companion regression of future short-rate changes on forward spreads, and Campbell--Shiller regressions for maturities of 2 to 10 years (Lab Figures~\ref{ch:10}.2 and~\ref{ch:10}.3). Standard errors\index{standard error} are Newey--West with 18 lags to account for the overlapping monthly observations.

\textbf{Part B, Cochrane--Piazzesi.} The average excess return on 2- to 5-year bonds regressed on the one-year yield and four forward rates; the loading of each maturity on the resulting factor; and out-of-sample forecasts from 1985 with an expanding window\index{walk-forward evaluation!expanding window} and an 11-month gap so that no training target overlaps the forecast date. Three predictor sets are compared with the historical mean\index{historical-mean benchmark}: the full forward curve, a single forward spread and the level of rates alone (Lab Figure~\ref{ch:10}.4).

\textbf{Part C, subperiods.} The Fama--Bliss slope for 5-year bonds, the Campbell--Shiller slope at the longest maturity and the Cochrane--Piazzesi $R^2$ in four periods: the start to 1984, 1985--2007, 2008--2015 (the zero lower bound) and 2016 onward.

\textbf{Part D, the investor's view.} Average excess returns by maturity, and returns after steep versus flat curves, with ``steep'' defined using only past information (Lab Figure~\ref{ch:10}.5).

Two caveats apply. The Fed's curve is a fitted curve, not a set of traded bonds, and the classic studies used discount bonds constructed from CRSP data, which are not free. Results will be similar in spirit but not identical. Ten-year yields begin in 1971, so the longest maturities use a shorter sample.

\textbf{What the simulated run shows.} The simulation adds a time-varying term premium to the expected path of short rates, so the EH is false by construction. The lab finds what it should: Fama--Bliss slopes of about 0.6 to 0.7, significantly above zero; Campbell--Shiller slopes negative and well below 1; higher excess returns after steep curves. The tests on data where the EH holds exactly recover slopes of 0 and 1. These results confirm that the code works. They say nothing about real markets. Run the lab on the real data to reach your own verdict.

\section{Stress tests}
\label{sec:10-stress}

\textbf{Subperiods.} The postwar sample contains very different monetary regimes: the inflation of the 1970s, the Volcker disinflation, the long decline in rates from the 1980s, the zero lower bound after 2008 and the rapid tightening of 2022--23. A term premium that varies with inflation uncertainty would be most visible in the first half. Part C shows whether the rejection of the EH is a product of one era.

\textbf{The zero lower bound.} From 2008 to 2015 the short rate was pinned near zero. Forward rates still moved, driven by expectations of when the Fed would lift rates and by the term premium, but realized short rates did not. Regressions over this period behave differently, and the long-standing evidence may not carry over to it.

\textbf{Maturities.} The EH is a statement about all maturities. If it fails at two years but nearly holds at ten, or the reverse, that is informative about where the term premium lives.

\textbf{Out of sample.}\index{in-sample versus out-of-sample|textbf} In-sample $R^2$ values for bond return forecasts are inflated by persistent predictors and overlapping returns. The lab's out-of-sample test, which uses only data available at each date and leaves an 11-month gap, is the honest test. A negative out-of-sample $R^2$\index{out-of-sample $R^2$} would mean that an investor would have done better using the historical mean.

\textbf{Small samples.} Bauer and Hamilton (2018) showed that standard errors can understate uncertainty badly in these regressions. Their bootstrap\index{bootstrap} tests are beyond this lab, but Exercise 4 asks how far the lab's $t$-statistics can be trusted. Treat $t$-statistics near 2 with suspicion.

\section{Verdict}
\label{sec:10-verdict}

\begin{callout}
\textbf{Verdict: fails} (as of October 2026). Forward rates are not unbiased forecasts of future short rates. Bond risk premia vary over time, and the slope of the yield curve forecasts excess returns on long bonds, not just changes in short rates. The more ambitious claim, that a richer function of the forward curve forecasts bond returns with high $R^2$, is weaker than it once appeared, but that does not rescue the expectations hypothesis.
\end{callout}

What fails:

\begin{itemize}
  \item \textbf{Forward spreads forecast excess returns.} In the classic studies and their updates, Fama--Bliss slopes are positive and significant. The forward spread is partly a forecast of rates and partly a premium for holding long bonds.
  \item \textbf{Long yields move the wrong way.} Campbell--Shiller slopes are below 1, and often negative, across maturities. On average, a steep curve has been followed by falling, not rising, long yields.
  \item \textbf{The premium is counter-cyclical.} Expected bond excess returns are higher in bad times, which the EH rules out.
\end{itemize}

What survives:

\begin{itemize}
  \item \textbf{Forward rates contain information about future short rates}, especially at horizons of several years, because rates mean-revert. They are biased forecasts, not useless ones.
  \item \textbf{The term premium is hard to forecast precisely.} Small-sample corrections and out-of-sample tests weaken the case for large, exploitable predictability (Bauer and Hamilton 2018; Thornton and Valente 2012). The evidence is strong that the premium varies, and weaker on how well anyone can track it in real time.
\end{itemize}

The verdict would move to \emph{weakens} if the rejection of the EH were confined to one monetary regime and disappeared in the others.

\section{What investors can use}
\label{sec:10-investors}

\textbf{Read forward rates as forecasts plus a premium.} A market-implied path for policy rates, extracted from forwards or futures, includes a term premium. It overstates expected rate rises when the premium is positive and understates them when the premium is negative. Estimated term premiums, such as those published by the Federal Reserve Bank of New York from the Adrian, Crump and Moench (2013) model, are one way to separate the two, though they are model estimates with wide uncertainty.

\textbf{Carry and roll-down.} A bond investor in a steep curve earns two things if rates do not change: the extra yield of the longer bond (\textbf{carry}) and the price gain as the bond rolls down to lower yields at shorter maturities (\textbf{roll-down}\index{roll-down|textbf}). The EH says these are expected to be offset by rising rates. The evidence says they are not fully offset on average. Carry and roll-down are thus a reasonable guide to expected excess returns, not a guaranteed one.

\textbf{Duration is a bet.}\index{duration|textbf} Extending maturity when the curve is steep has been rewarded on average, but with the volatility\index{volatility} of long bonds. In 2022, when rates rose sharply, long Treasuries suffered some of their worst losses on record.

\textbf{Questions to ask a bond manager:}

\begin{itemize}
  \item How much of the portfolio's expected return comes from carry and roll-down, and what happens if rates rise by the forward rates?
  \item Does the strategy assume forward rates are forecasts of future rates, or does it allow for a term premium?
  \item How did the strategy perform when the curve was inverted\index{yield curve!inverted}, and in 2022?
\end{itemize}

\begin{callout}
\textbf{What Failed: reading the yield curve as a pure forecast of policy.} Between June 2004 and June 2006 the Federal Reserve raised its policy rate from 1\% to 5.25\%. Under the expectations hypothesis, long-term yields should have risen as the path of short rates moved up. Instead, long yields fell during the first year of the tightening and ended the period little higher than where they began. In February 2005, Chairman Alan Greenspan called the behaviour of long rates a ``conundrum'' (Greenspan 2005).

The conundrum was a puzzle only if long yields were pure forecasts of short rates. Term-premium models attributed much of the decline to a falling term premium, which various explanations linked to lower inflation uncertainty, foreign official purchases of Treasuries and demand from pension funds. Investors and policymakers who read the long end of the curve as the market's forecast of policy rates misread both the market and its signal. The same error recurs whenever a change in the term premium is taken for a change in expected policy.
\end{callout}

\section{Open questions}
\label{sec:10-open}

\begin{itemize}
  \item \textbf{Has central-bank bond buying changed the term premium?} Large-scale asset purchases after 2008 and in 2020 were designed to lower long yields by reducing the term premium. Did they change its behaviour or its predictability?
  \item \textbf{Is the Cochrane--Piazzesi factor real?} Robust small-sample tests weaken it. Does it survive in other countries and in data after 2003?
  \item \textbf{What drives the term premium?} Inflation uncertainty, risk aversion, supply and demand all play a part. Which dominates now?
  \item \textbf{Do macroeconomic variables add to the yield curve?} Several studies find that measures of real activity forecast bond returns beyond yields. Does that evidence survive the same robust tests?
\end{itemize}

\section*{Exercises}

\subsection*{Conceptual (all readers)}

\begin{enumerate}
  \item The 1-, 2- and 3-year zero yields are 3\%, 3.5\% and 4\% (continuously compounded). Compute the forward rates $f(2)$ and $f(3)$. Under the expectations hypothesis, what one-year rates are expected one and two years from now?
  \item Using the yields in Exercise 1, compute the one-year excess return on the 3-year bond if the yield curve is unchanged a year from now. What would the excess return be if the expectations hypothesis held?
  \item Show that the Fama--Bliss slopes for excess returns and for short-rate changes add up to 1 when the horizon is one year ($n = 2$).
\end{enumerate}

\subsection*{Lab extensions (technical readers)}

\begin{enumerate}
  \setcounter{enumi}{3}
  \item Part A: at which maturities is the Fama--Bliss slope significantly above 0? Is any significantly different from 1? Given Bauer and Hamilton's (2018) warning about small-sample distortions, which of these conclusions would you still trust, and why?
  \item Part A: how does the Campbell--Shiller slope change with maturity? Explain in words what a negative slope means for a bond investor.
  \item Part B: compare the in-sample and out-of-sample $R^2$ of the Cochrane--Piazzesi regression. Which of the three predictor sets does best out of sample?
  \item Part C: in which subperiod is the evidence against the hypothesis strongest? What happened in 2008--2015?
  \item Part D: build a strategy that holds the 5-year bond after a steep curve and the 1-year bond otherwise. Compare its average excess return and Sharpe ratio with always holding the 5-year bond.
  \item A newspaper reports that ``markets expect three rate cuts next year,'' based on futures prices. Using this chapter's evidence, explain what the futures prices actually tell us.
\end{enumerate}
